\section{Methodology}
\label{sec:method}

\subsection{Research Question and Experimental Design}
We test whether self-critique induces structured internal changes rather than only output variation. For each question, we generate one \direct draft and up to four revisions: \rewrite control, \selfone, \selfthree, and \external. We then compare both answers and hidden representations between pre- and post-revision states.

\subsection{Datasets and Splits}
We use local copies of two reasoning datasets: \gsm test split (1,319 items) and \csqa validation split (1,221 items). The main mechanistic run uses a deterministic subset of 25 \gsm + 25 \csqa examples ($n=50$). A size-comparison run uses 12 + 12 examples ($n=24$). A behavioral-only API run uses 20 sampled items.

Data quality checks reported zero missing values, zero detected duplicates, and complete gold-label parse success for \gsm (1319/1319). For \csqa, class counts are near balanced (A:239, B:255, C:241, D:251, E:235).

\subsection{Models and Implementation}
Local mechanistic runs use \qwenonefive (primary) and \qwenzerofive (size comparison). Behavioral triangulation uses \gptfourone via API. We run deterministic decoding with temperature 0 and max 48 new tokens in the main run. Hidden states are extracted with \texttt{output\_hidden\_states=True}, using final-token residual vectors from depth-spanning selected layers (for \qwenonefive: layers 1, 7, 14, 21, 28).

Environment: Python 3.12.8, PyTorch 2.10.0+cu128, Transformers 5.3.0, Datasets 4.4.0, NumPy 2.3.5, SciPy 1.17.1, scikit-learn 1.8.0, statsmodels 0.14.6. Hardware includes two RTX 3090 GPUs (24GB each) with bfloat16 inference.

\subsection{Metrics}
We evaluate:
\begin{itemize}
    \item \textbf{Behavior}: exact-match accuracy for \gsm numeric answers and \csqa option labels.
    \item \textbf{Drift}: cosine distance between draft and revised activations.
    \item \textbf{Geometry retention}: centered kernel alignment (CKA) between pre/post representation matrices.
    \item \textbf{Correctness separability}: logistic-probe AUC/F1 for correct vs. incorrect outcomes.
    \item \textbf{Trajectory stability}: round-to-round drift across three self-critique iterations.
\end{itemize}

\subsection{Statistical Analysis and Baselines}
The baseline is \direct prompting. We compare each revision condition to \direct with paired Wilcoxon signed-rank tests and apply Benjamini--Hochberg FDR correction across comparisons. Probe results are reported descriptively due limited sample size, especially for the small-model run where class imbalance makes AUC/F1 unstable.
